\documentclass[10pt,a4paper]{amsart}
\usepackage{amsmath,amssymb,amsthm,graphicx}
% For commutative diagrams we use Xy-pic
\usepackage[all]{xy}

% We use 2cell for 2-commutative diagrams.
\xyoption{2cell}
% Theorem environments.
%
\theoremstyle{plain}
\newtheorem{theorem}[subsection]{Theorem}
\newtheorem{proposition}[subsection]{Proposition}
\newtheorem{lemma}[subsection]{Lemma}

\theoremstyle{definition}
\newtheorem{definition}[subsection]{Definition}
\newtheorem{example}[subsection]{Example}
\newtheorem{exercise}[subsection]{Exercise}
\newtheorem{situation}[subsection]{Situation}

\theoremstyle{remark}
\newtheorem{remark}[subsection]{Remark}
\newtheorem{remarks}[subsection]{Remarks}

\numberwithin{equation}{subsection}

% Macros
%
\def\lim{\mathop{\mathrm{lim}}\nolimits}
\def\colim{\mathop{\mathrm{colim}}\nolimits}
\def\Spec{\mathop{\mathrm{Spec}}}
\def\Hom{\mathop{\mathrm{Hom}}\nolimits}
\def\Ext{\mathop{\mathrm{Ext}}\nolimits}
\def\SheafHom{\mathop{\mathcal{H}\!\mathit{om}}\nolimits}
\def\SheafExt{\mathop{\mathcal{E}\!\mathit{xt}}\nolimits}
\def\Sch{\mathit{Sch}}
\def\Mor{\mathop{\mathrm{Mor}}\nolimits}
\def\Ob{\mathop{\mathrm{Ob}}\nolimits}
\def\Sh{\mathop{\mathit{Sh}}\nolimits}
\def\NL{\mathop{N\!L}\nolimits}
\def\CH{\mathop{\mathrm{CH}}\nolimits}
\def\proetale{{pro\text{-}\acute{e}tale}}
\def\etale{{\acute{e}tale}}
\def\QCoh{\mathit{QCoh}}
\def\Ker{\mathop{\mathrm{Ker}}}
\def\Im{\mathop{\mathrm{Im}}}
\def\Coker{\mathop{\mathrm{Coker}}}
\def\Coim{\mathop{\mathrm{Coim}}}

% Boxtimes
%
\DeclareMathSymbol{\boxtimes}{\mathbin}{AMSa}{"02}

%
% Macros for moduli stacks/spaces
%
\def\QCohstack{\mathcal{QC}\!\mathit{oh}}
\def\Cohstack{\mathcal{C}\!\mathit{oh}}
\def\Spacesstack{\mathcal{S}\!\mathit{paces}}
\def\Quotfunctor{\mathrm{Quot}}
\def\Hilbfunctor{\mathrm{Hilb}}
\def\Curvesstack{\mathcal{C}\!\mathit{urves}}
\def\Polarizedstack{\mathcal{P}\!\mathit{olarized}}
\def\Complexesstack{\mathcal{C}\!\mathit{omplexes}}
% \Pic is the operator that assigns to X its picard group, usage \Pic(X)
% \Picardstack_{X/B} denotes the Picard stack of X over B
% \Picardfunctor_{X/B} denotes the Picard functor of X over B
\def\Pic{\mathop{\mathrm{Pic}}\nolimits}
\def\Picardstack{\mathcal{P}\!\mathit{ic}}
\def\Picardfunctor{\mathrm{Pic}}
\def\Deformationcategory{\mathcal{D}\!\mathit{ef}}
\usepackage{cleveref}
\crefname{originallemma}{original lemma}{original lemmas}
\Crefname{originallemma}{Original lemma}{Original lemmas}
\makeatletter
\newlabel{more-algebra-lemma-divide-by-torsion}{{05EB}{source}{}{sourcepointer-1}{}}
\newlabel{more-algebra-lemma-divide-by-torsion@cref}{{[originallemma][0][]05EB}{[1][0][]source}{}{}{}}
\makeatother
\crefname{originallemma}{original lemma}{original lemmas}
\Crefname{originallemma}{Original lemma}{Original lemmas}
\makeatletter
\newlabel{lemma-compare-periodic-lengths}{{0EA9}{source}{}{sourcepointer-2}{}}
\newlabel{lemma-compare-periodic-lengths@cref}{{[originallemma][0][]0EA9}{[1][0][]source}{}{}{}}
\makeatother
\crefname{originallemma}{original lemma}{original lemmas}
\Crefname{originallemma}{Original lemma}{Original lemmas}
\makeatletter
\newlabel{lemma-additivity-periodic-length}{{0EA7}{source}{}{sourcepointer-3}{}}
\newlabel{lemma-additivity-periodic-length@cref}{{[originallemma][0][]0EA7}{[1][0][]source}{}{}{}}
\makeatother
\crefname{originalequation}{original equation}{original equations}
\Crefname{originalequation}{Original equation}{Original equations}
\makeatletter
\newlabel{equation-multiplicity-coker-ker}{{0EA6}{source}{}{sourcepointer-4}{}}
\newlabel{equation-multiplicity-coker-ker@cref}{{[originalequation][0][]0EA6}{[1][0][]source}{}{}{}}
\makeatother
\title{Multiplying one structure map of a two-one-periodic complex by a parameter of vanishing support}
\author{The Stacks Project Authors}
\date{stacks-project:chow.tex:local-snapshot-20261011; GNU FDL 1.2 or later}
\begin{document}
\maketitle
\noindent\textit{Editorial note.} All delivered mathematical text is the author's own, copied from the pinned Stacks Project chapter file chow.tex; no statement or proof step was rewritten, reordered or removed. The author proves only the first of the two displayed formulas and says so; this is reproduced as written.\par\medskip
\noindent\textit{Editorial note.} The preamble is the chapter's own native preamble excerpt used for this source: the Xy-pic setup at lines 14-18 and the theorem declarations with the mathematical macros at lines 154-217 of sources/stacks-chow/preamble.tex. The quoted class line and the xr-hyper externaldocument block are not included because the assembler supplies a plain amsart wrapper; no macro was invented or redefined.\par\medskip
\noindent\textit{Editorial note.} Same-chapter references are located at their \textbackslash{}label lines inside the pinned chapter file; the cross-chapter pointer is located in the project's own tag list staged in the snapshot (sources/stacks-chow/tags.txt), at the verbatim line carrying its canonical tag.\par\medskip
\noindent\textit{Allowed standard prerequisites (editorial).}\begin{itemize}
\item Two-one-periodic complexes and their multiplicity e\_R. Same chapter, Definition environment labelled definition-periodic-complex (source line 110) together with the displayed definition of e\_R at source lines 200-205 (label equation-multiplicity-coker-ker, line 203). Supplies the objects (M, phi, psi), the quantity e\_R and the displayed identity used at proof line 509; the two formulas are the counted content.
\item x-power torsion submodule and x-power torsion freeness of the quotient. More on Algebra, Lemma labelled more-algebra-lemma-divide-by-torsion (canonical Stacks tag 05EB); invoked at proof lines 483-484. Supplies that M'' = M/M[x\textasciicircum{}infinity] is x-power torsion free; the reduction to that quotient is retained in the counted span.
\item Vanishing criteria for the multiplicities of the summands. Same chapter, Lemma labelled lemma-compare-periodic-lengths (source line 283, Stacks tag 0EA9), invoked at proof lines 494-496; and Lemma labelled lemma-additivity-periodic-length (source line 209, Stacks tag 0EA7), invoked at proof line 497. Supply the three vanishing equalities for the torsion part and the additivity used to discard it; the reduction and the final exact sequence are inside the counted span.
\end{itemize}\medskip
\ifdefined\hypertarget\hypertarget{sourcepointer-1}{}\fi
\noindent\textit{Source reference.} Original lemma 05EB: sources/stacks-chow/tags.txt, source lines 6984--6984. Cross-chapter standard prerequisite invoked once at proof lines 483-484 for the x-power torsion freeness of the quotient; the reduction built on it is retained in the counted span.\par\medskip
\ifdefined\hypertarget\hypertarget{sourcepointer-2}{}\fi
\noindent\textit{Source reference.} Original lemma 0EA9: sources/stacks-chow/pinned-original-source.tex, source lines 283--283. Same-chapter pointer invoked at proof lines 494-496 for the three vanishing equalities on the torsion part; the reduction and the final computation are retained in the counted span.\par\medskip
\ifdefined\hypertarget\hypertarget{sourcepointer-3}{}\fi
\noindent\textit{Source reference.} Original lemma 0EA7: sources/stacks-chow/pinned-original-source.tex, source lines 209--209. Same-chapter pointer invoked at proof line 497 for additivity of e\_R, which is what allows discarding the torsion part.\par\medskip
\ifdefined\hypertarget\hypertarget{sourcepointer-4}{}\fi
\noindent\textit{Source reference.} Original equation 0EA6: sources/stacks-chow/pinned-original-source.tex, source lines 203--203. Same-chapter pointer to the displayed definition of e\_R, invoked at proof line 509 to convert the exact sequence into the asserted identity; the exact sequence itself is retained in the counted span.\par\medskip
\begin{lemma}
\label{lemma-multiply-period-length}
Let $(R, \mathfrak m)$ be a Noetherian local ring. Let
$(M, \varphi, \psi)$ be a $(2, 1)$-periodic complex over $R$
with $M$ finite and with cohomology groups of finite length over $R$.
Let $x \in R$ be such that $\dim(\text{Supp}(M/xM)) \leq 0$. Then
$$
e_R(M, x\varphi, \psi) = e_R(M, \varphi, \psi) - e_R(\Im(\varphi), 0, x)
$$
and
$$
e_R(M, \varphi, x\psi) = e_R(M, \varphi, \psi) + e_R(\Im(\psi), 0, x)
$$
\end{lemma}

\begin{proof}
We will only prove the first formula as the second is proved
in exactly the same manner.
Let $M' = M[x^\infty]$ be the $x$-power torsion submodule of $M$.
Consider the short exact sequence $0 \to M' \to M \to M'' \to 0$.
Then $M''$ is $x$-power torsion free (More on Algebra, Lemma
\ref{more-algebra-lemma-divide-by-torsion}).
Since $\varphi$, $\psi$ map $M'$ into $M'$
we obtain a short exact sequence
$$
0 \to (M', \varphi', \psi') \to (M, \varphi, \psi) \to
(M'', \varphi'', \psi'') \to 0
$$
of $(2, 1)$-periodic complexes. Also, we get a short exact sequence
$0 \to M' \cap \Im(\varphi) \to \Im(\varphi) \to \Im(\varphi'') \to 0$.
We have
$e_R(M', \varphi, \psi) = e_R(M', x\varphi, \psi) =
e_R(M' \cap \Im(\varphi), 0, x) = 0$
by Lemma \ref{lemma-compare-periodic-lengths}.
By additivity (Lemma \ref{lemma-additivity-periodic-length})
we see that it suffices to prove the lemma for $(M'', \varphi'', \psi'')$.
This reduces us to the case discussed in the next paragraph.

\medskip\noindent
Assume $x : M \to M$ is injective.
In this case $\Ker(x\varphi) = \Ker(\varphi)$.
On the other hand we have a short exact sequence
$$
0 \to \Im(\varphi)/x\Im(\varphi) \to
\Ker(\psi)/\Im(x\varphi) \to \Ker(\psi)/\Im(\varphi) \to 0
$$
This together with (\ref{equation-multiplicity-coker-ker}) proves the formula.
\end{proof}

\end{document}
